\appendix
\section{Appendix}
\label{sec:appendix}

% Mirrors slides.qmd "Appendix" outline. One subsection per outline item.

The two algorithms below are used throughout the notes. They are stated here in
full so the main text can refer to them without interrupting its argument.

\subsection{TT-Cross}
\label{subsec:appendix-tt-cross}
% Keywords: tensor cross interpolation algorithm details, maxvol/rectangular
% maxvol, skeleton decomposition, pivoting strategy, convergence and error
% bounds, computational cost.

TT-cross builds a tensor train for a black-box $f$ from $O(n\chi^2)$
evaluations, never forming the full array~\cite{oseledets2010ttcross}. It is
what makes \cref{subsec:finding-qtt,subsec:tt-cross-geometry} possible.

\paragraph{The matrix case first.}
\begin{itemize}
  \item A rank-$r$ matrix $M$ is exactly reconstructed from $r$ rows $I$ and $r$
    columns $J$ by the \textbf{skeleton} (or CUR) formula
    \begin{equation}
      M = M(:,J)\, M(I,J)^{-1} M(I,:) .
      \label{eq:skeleton}
    \end{equation}
  \item The whole question is which rows and columns to pick. A badly
    conditioned $M(I,J)$ destroys the reconstruction, since its inverse appears
    in~\eqref{eq:skeleton}.
  \item \textbf{maxvol} answers it: choose the $r \times r$ submatrix of maximum
    volume, i.e.\ largest $\abs{\det}$.\margintext{The maximum-volume principle
    is due to Goreinov and Tyrtyshnikov (2001); Oseledets and Tyrtyshnikov
    carried it up to tensor trains in 2010.} For an approximately rank-$r$ $M$
    the resulting error is within a factor $(r+1)$ of the best rank-$r$
    approximation in the Chebyshev norm.
  \item Finding the true maximum-volume submatrix is NP-hard. The practical
    algorithm swaps rows greedily until no single swap increases $\abs{\det}$,
    which lands within a controlled factor of the true maximum and costs
    $O(N r^2)$ per sweep.
\end{itemize}

\paragraph{The tensor-train algorithm.}
Maintain, for each bond $k$, a set of row multi-indices $\mathcal{I}_k$ and
column multi-indices $\mathcal{J}_k$. Then:

\begin{enumerate}
  \item Initialize $\mathcal{I}_k$ and $\mathcal{J}_k$, either at random or from
    a coarse pass over the domain.
  \item Sweeping left to right, form the slice
    $F_k = f(\mathcal{I}_k,\, i_k,\, \mathcal{J}_k)$, a $\chi \times d\chi$
    matrix obtained by varying only the local index $i_k$ while the pivots hold
    everything else fixed. This is the only place $f$ is called.
  \item Apply \textbf{maxvol} to $F_k$: select the $\chi$ rows of maximum
    $\abs{\det}$, and use them to update $\mathcal{I}_{k+1}$.
  \item Set the core to $F_k \hat{F}_k^{-1}$, where $\hat{F}_k$ is the selected
    maximum-volume submatrix, and carry the residual factor to the right.
  \item Sweep back right to left and repeat until the cross error, measured on
    fresh sample points, plateaus.
\end{enumerate}

\begin{remark}[Interpolation property]
  The resulting MPS is \textbf{exact} on every fibre it sampled --- that is what
  $\hat{F}_k^{-1}$ in step~4 enforces. This is the algorithm's defining
  strength and its defining weakness: it reproduces sampled values exactly,
  including sampled noise, and it says nothing at all about the fibres it never
  looked at.
\end{remark}

\paragraph{Practical notes.}
\begin{itemize}
  \item Cost is $O(n \chi^3)$ arithmetic and $O(n \chi^2)$ function evaluations
    for $n$ tensors of bond dimension $\chi$ --- \textbf{independent of $2^n$}.
  \item \textbf{Rectangular maxvol} keeps more than $\chi$ pivots, improving
    robustness at modest extra cost. Generally worth switching on.
  \item \textbf{Rank-adaptive variants} insert pivots where the local error is
    largest, rather than fixing $\chi$ in advance~\cite{nunez2025tci}.
  \item Convergence is heuristic. There is no guarantee for adversarial or noisy
    black boxes, and the failure modes of \cref{subsec:finding-qtt} are all
    observed in practice. Always validate on points the cross never sampled.
\end{itemize}

\subsection{DMRG}
\label{subsec:appendix-dmrg}
% Keywords: density matrix renormalization group, two-site sweep algorithm,
% local optimization/eigensolver, sweeping schedule, connection to
% variational MPS optimization, adaptive bond-dimension growth.

DMRG~\cite{white1992,schollwoeck2011} is the variational optimization of an MPS
by local updates. In these notes it appears as the linear solver of
\cref{subsec:linear-solvers}, as the pressure-Poisson step of
\cref{subsec:qtt-cfd-frameworks}, and as the all-at-once space--time solve of
\cref{subsec:dmd-space-time}.

\paragraph{The two-site sweep.}
\begin{enumerate}
  \item Bring the MPS into mixed canonical form with the orthogonality centre at
    site $k$ (\cref{subsec:mps}).
  \item Merge sites $k$ and $k+1$ into a single two-site tensor of size
    $\chi d^2 \chi$.
  \item Solve the \textbf{local problem} in that tensor alone, with every other
    tensor held fixed: an eigenproblem for ground states, a linear system for
    $Ax = b$. The effective operator is assembled from cached left and right
    \textbf{environments} --- the partial contractions of the MPO with the
    frozen parts of the chain --- which are updated one site at a time rather
    than rebuilt.
  \item Split the result by SVD, truncating to a tolerance $\varepsilon$ or a
    rank cap $\chi_{\max}$.
  \item Move the orthogonality centre one site along and repeat, sweeping right
    then left.
\end{enumerate}

\begin{remark}
  Because the canonical form makes the environments orthonormal, the local
  problem in step~3 is a genuine restriction of the global one, not an
  approximation of it. This is why the global objective decreases monotonically.
\end{remark}

\paragraph{Cost and practical notes.}
\begin{itemize}
  \item Per sweep, $O(n \chi^3 d^2 + n \chi^2 d\, \chi_{\mathrm{MPO}})$ for the
    contractions, with the local solve reaching $O(\chi^6)$ in the worst case.
    Observed totals are much better (\cref{subsec:linear-solvers}).
  \item \textbf{Adaptive bond dimension} is the point of the two-site variant:
    the SVD in step~4 may return a larger rank than it started with, so
    $\varepsilon$ rather than $\chi$ becomes the control parameter.
  \item \textbf{Escaping local minima.} Two-site updates explore directions
    unavailable to any single tensor, so they escape stalls that trap one-site
    sweeps. AMEn recovers most of this at one-site cost by enriching the update
    with residual directions~\cite{dolgov2014amen,holtz2012als}.
  \item \textbf{Sweeping schedule.} Start with a loose tolerance and a low rank
    cap, then tighten both over successive sweeps. Two or three sweeps usually
    suffice from a warm start.
  \item \textbf{Warm starts} matter more than any other tuning parameter in a
    time-stepping context, because consecutive solutions are close.
\end{itemize}
